\documentclass[11pt,a4paper]{article}
\usepackage{fontspec,amsmath,amssymb,amsthm,microtype}
\setmainfont{lmroman10-regular.otf}[BoldFont=lmroman10-bold.otf,ItalicFont=lmroman10-italic.otf,BoldItalicFont=lmroman10-bolditalic.otf]
\setmonofont{lmmono10-regular.otf}[Scale=MatchLowercase]
\usepackage[margin=25mm]{geometry}
\usepackage[unicode,hidelinks]{hyperref}
\setlength{\emergencystretch}{3em}
\setlength{\parskip}{0.4em}
\setlength{\parindent}{0pt}
\allowdisplaybreaks[2]
\raggedbottom
\title{\texorpdfstring{Assemble continuous Brownian reference segments}{Assemble continuous Brownian reference segments}}
\author{LangevinLean mathematical exposition}
\date{}
\begin{document}\maketitle
\noindent Source module: \texttt{Unadjusted\allowbreak{}Langevin/\allowbreak{}Brownian\allowbreak{}Reference.\allowbreak{}lean}. The accompanying HTML pairs this mathematical exposition with the exact Lean statements and records the scope of the formalization.


Construct consecutive \(L^2\) Langevin segments from actual Brownian-coordinate noise and package them as an analytic reference. At this intermediate stage Gibbs stationarity remains an explicit premise, later discharged in CanonicalULA.


\section{\texorpdfstring{brownian\allowbreak{}Segment\allowbreak{}From}{brownianSegmentFrom}}


Choose the integral solution supplied by \texttt{exists\_\allowbreak{}lp\_\allowbreak{}langevin\_\allowbreak{}integral\_\allowbreak{}of\_\allowbreak{}brownian} for the segment starting at absolute time \(nh\), with the supplied initial \(L^2\) variable. Classical choice selects data from a proved existence theorem; its defining equation and continuity are recovered by the specification theorem.



{\raggedright\noindent Lean declarations: \texttt{brownian\allowbreak{}Segment\allowbreak{}From}.\par}

\section{\texorpdfstring{brownian\allowbreak{}Segment\allowbreak{}From spec}{brownianSegmentFrom spec}}


Apply the choice specification to obtain continuity, the time-zero initial value, the gradient integral equation on \([0,h]\), and the exact noise norm. These four properties describe the selected segment independently of its noncomputable implementation.



{\raggedright\noindent Lean declarations: \texttt{brownian\allowbreak{}Segment\allowbreak{}From\_\allowbreak{}spec}.\par}

\section{\texorpdfstring{brownian\allowbreak{}Knots}{brownianKnots}}


Set the zeroth knot to the prescribed initial variable. Define the next knot to be the endpoint at local time \(h\) of the segment started at the current knot. This recursion ensures consecutive segments meet exactly as \(L^2\) classes.



{\raggedright\noindent Lean declarations: \texttt{brownian\allowbreak{}Knots}.\par}

\section{\texorpdfstring{brownian\allowbreak{}Segments}{brownianSegments}}


At index \(n\), select \texttt{brownian\allowbreak{}Segment\allowbreak{}From} with the recursively constructed \(n\)-th knot as initial state. The absolute noise start is \(nh\), so all segments are driven by increments of the same Brownian process.



{\raggedright\noindent Lean declarations: \texttt{brownian\allowbreak{}Segments}.\par}

\section{\texorpdfstring{brownian\allowbreak{}Segments start}{brownianSegments start}}


The chosen-segment specification states its value at zero equals its supplied initial state. Here that state is the \(n\)-th knot, so \(Y^{(n)}_0=K_n\) follows directly by specialization.



{\raggedright\noindent Lean declarations: \texttt{brownian\allowbreak{}Segments\_\allowbreak{}start}.\par}

\section{\texorpdfstring{brownian\allowbreak{}Segments initial}{brownianSegments initial}}


Apply the segment-start identity at \(n=0\), then unfold the zeroth knot. The actual constructed reference therefore begins at exactly the prescribed random variable, preserving its chosen joint realization with the numerical initialization.



{\raggedright\noindent Lean declarations: \texttt{brownian\allowbreak{}Segments\_\allowbreak{}initial}.\par}

\section{\texorpdfstring{brownian\allowbreak{}Euler\allowbreak{}Noise}{brownianEulerNoise}}


Normalize the actual segment endpoint noise by \(1/\sqrt{2h}\). This defines the \(L^2\) innovation used by the Euler chain, so it shares the exact Brownian increment driving the reference segment.



{\raggedright\noindent Lean declarations: \texttt{brownian\allowbreak{}Euler\allowbreak{}Noise}.\par}

\section{\texorpdfstring{brownian\allowbreak{}Euler\allowbreak{}Noise endpoint}{brownianEulerNoise endpoint}}


For \(h>0\), \(\sqrt{2h}\ne0\). Multiply the normalized noise class by \(\sqrt{2h}\) and cancel the inverse scalar. The resulting equality identifies the Euler noise contribution with the same Brownian endpoint input in the integral equation.



{\raggedright\noindent Lean declarations: \texttt{brownian\allowbreak{}Euler\allowbreak{}Noise\_\allowbreak{}endpoint}.\par}

\section{\texorpdfstring{brownian\allowbreak{}Reference}{brownianReference}}


Populate \texttt{L2\allowbreak{}Reference} with the constructed segments. Their specification proves continuity and the integral equation; recursive knots prove endpoint compatibility; normalized noise proves endpoint matching; Gibbs-score moment transfer proves the drift bound. The supplied all-time Gibbs-law premise remains visible and is later discharged by the stationarity construction.



{\raggedright\noindent Lean declarations: \texttt{brownian\allowbreak{}Reference}.\par}

\section{\texorpdfstring{brownian\allowbreak{}Reference of strong\allowbreak{}Convex}{brownianReference of strongConvex}}


Use strong convexity to derive finiteness of the Gibbs partition, then call \texttt{brownian\allowbreak{}Reference}. All analytic, normalization, moment, and noise fields are proved internally. The all-time law premise is still explicit in this constructor and must be provided by the stationary-law theorem.



{\raggedright\noindent Lean declarations: \texttt{brownian\allowbreak{}Reference\_\allowbreak{}of\_\allowbreak{}strong\allowbreak{}Convex}.\par}

\end{document}
